% Solution_Template.tex
\documentclass[11pt,a4paper]{article}
\usepackage[utf8]{inputenc}
\usepackage[T1]{fontenc}
\usepackage[ngerman, english]{babel} % remove ngerman if you prefer English only
\usepackage{geometry}
\usepackage{fancyhdr}
\usepackage{hyperref}
\usepackage{array,etoolbox}
\usepackage{afterpage}
\usepackage{longtable}
\usepackage{bigdelim}
\usepackage{multirow}


\preto\tabular{\setcounter{magicrownumbers}{-1}}
\newcounter{magicrownumbers}
\newcommand\rownumber{\stepcounter{magicrownumbers}\arabic{magicrownumbers}}

\geometry{margin=25mm}
\pagestyle{fancy}
\fancyhf{}
\renewcommand{\headrulewidth}{0.4pt}
\lhead{\textbf{Name:} \studentname}
\rhead{\textbf{Matr.-Nr.:} \studentnumber}
\cfoot{\thepage}

% User info - edit these
\newcommand{\studentname}{Matthias Janßen}
\newcommand{\studentnumber}{1871808}
\newcommand{\course}{Übersetzung und Analyse von Programmen}
\newcommand{\institute}{Universität Trier}

\title{\course}
\author{\studentname\\Matrikelnummer: \studentnumber\\\institute}
\date{\today}

\begin{document}

\maketitle

\begin{center}
    \textbf{Abgabe:} Erstes Übungsblatt \\
    \vspace{4mm}
    \textbf{Betreuer:} Prof. Dr. Stephan Diel \\ Jan Niclas Ruppenthal, M.Sc.
\end{center}

\bigskip

\section*{Aufgaben}

\section{Aufgabe 1: Programm der MaMa}
\begin{verbatim}
    1) Sei L = [n1, n2]
    2) Falls n1 = 0 gehe zu 5
    3) Falls n1 > n2 dann n1 - n2 sonst n2 - n1
    4) Falls n2 != 0 gehe zu 3
    4) Fertig n1 ist der ggT
    5) Fertig n2 ist der ggT

\end{verbatim}
\begin{itemize}
    \item P(0) = ldo(-1)
    \item p(1) = push(0)
    \item p(2) = equal(19)
    \item p(3) = ldo(0)
    \item p(4) = push(0)
    \item p(5) = equal(21)
    \item p(6) = ldo(-1)
    \item p(7) = ldo(-1)
    \item p(8) = leq(14)
    \item (p9) = ldo(-1)
    \item p(10) = ldo(-1)
    \item p(11) = sub
    \item p(12) = sto(-2)
    \item p(13) = ujp(3)
    \item p(14) = ldo(0)
    \item p(15) = ldo(-2)
    \item p(16) = sub
    \item p(17) = sto(-1)
    \item p(18) = ujp(3)
    \item p(19) = sto(-1)
    \item p(20) = stop
    \item p(21) = pop
    \item p(22) = stop
\end{itemize}


\section{Aufgabe 2: Rechnung der MaMa }
%\preto\tabular{\setcounter{magicrownumbers}{0}}


\begin{longtable}{@{\makebox[3em][r]{\rownumber\space}} c c c}
    (1,0, [4, 6]) & & \\
    (2, 1, [4, 6, 4]) &   ldo(-1)  & \rdelim\}{3}{1.8cm}[\normalfont a = 0 check]\\
    (3, 2, [4, 6, 4, 0]) & push(0)  & \\
    (1, 3, [4, 6]) & equal(19)  & \\
    (2, 4, [4, 6, 6]) & ldo(0)  & \rdelim\}{11}{1.8cm}[\normalfont 1. Iteration] \\
    (3, 5, [4, 6, 6, 0]) & push(0) & \\
    (1, 6, [4, 6]) & equal(21)  & \\
    (2, 7, [4, 6, 4]) & ldo(-1) & \\
    (3, 8, [4, 6, 4, 6]) & ldo(-1) & \\
    (1, 14, [4, 6]) & leq(14)  & \\
    (2, 15, [4, 6, 6]) & ldo(0)  & \\
    (3, 16, [4, 6, 6, 4]) & ldo(-2)  & \\
    (2, 17, [4, 6, 2]) & sub  & \\
    (1, 18, [4, 2]) & sto(-1)  & \\
    (1, 3, [4, 2]) & ujp(3)  & \\
    (2, 4, [4, 2, 2]) & ldo(0)  & \rdelim\}{11}{1.8cm}[\normalfont 2. Iteration]\\
    (3, 5, [4, 2, 2, 0]) & push(0)  & \\
    (1, 6, [4, 2]) & equal(21)  & \\
    (2, 7, [4, 2, 4]) & ldo(-1)  & \\
    (3, 8, [4, 2, 4, 2]) & ldo(-1)  & \\
    (1, 9, [4, 2]) & leq(14)  & \\
    (2, 10, [4, 2, 4]) & ldo(-1)  & \\
    (3, 11, [4, 2, 4, 2]) & ldo(-1)  & \\
    (2, 12, [4, 2, 2]) & sub  & \\
    (1, 13, [2, 2]) & sto(-2)  & \\
    (1, 3, [2, 2]) & ujp(3)  & \\
    (2, 4, [2, 2, 2]) & ldo(0)  & \rdelim\}{11}{1.8cm}[\normalfont 3. Iteration] \\
    (3, 5, [2, 2, 2, 0]) & push(0)  & \\
    (1, 6, [2, 2]) & equal(21)  & \\
    (2, 7, [2, 2, 2]) & ldo(-1)  & \\
    (3, 8, [2, 2, 2, 2]) & ldo(-1)  & \\
    (1, 14, [2, 2]) & leq(14)  & \\
    (2, 15, [2, 2, 2]) & ldo(0)  & \\
    (3, 16, [2, 2, 2, 2]) & ldo(-2)  & \\
    (2, 17, [2, 2, 0]) & sub  & \\
    (1, 18, [2, 0]) & sto(-1)  & \\
    (1, 3, [2, 0]) & ujp(3)  & \\
    (2, 4, [2, 0, 0]) & ldo(0)  & \rdelim\}{5}{1.8cm}[\normalfont Letzte Iteration] \\
    (3, 5, [2, 0, 0, 0]) & push(0)  & \\
    (1, 21, [2, 0]) & equal(21)  & \\
    (2, 22, [2]) & pop  & \\
    (0, 22, [2]) & stop  & \\

\end{longtable}

\bigskip Das Ergebnis ist 2, also der ggT von 4 und 6. Der Algorithmus hat eine Laufzeit von 41 Schritten
    für die gegebene Eingabe.

\section{Aufgabe 3: Laufzeit der MaMa}
    Laufzeizt: x * 11 + 8 \\
    Für die Laufzeit ist irrelevant ob n1 > n2 oder n2 > n1 gilt, da in beiden Fällen die gleiche Anzahl an Schritten benötigt wird.
    In jedem Schleifendurchlauf werden 11 Schritte ausgeführt, bis eine der beiden Zahlen 0 ist.
    Zusätzlich werden 8 Schritte für die Initialisierung und das Beenden des Programms benötigt.
    Der worst case tritt ein, wenn eine der beiden Zahlen 1 ist da dann bei n1 > n2 n1-1 Durchläufe benötigt werden.
\end{document}